\textbf{Pájky} se podle teploty tavení dělí na měkké a tvrdé. U měkkých pájek přednáška uvádí:
\begin{itemize}
\item \textbf{Olovnaté:} slitiny Sn--Pb, např. SnPb 63/37.
\item \textbf{Bezolovnaté:} Sn s přídavky Ag, Cu, Zn, Bi či Sb, případně Ge; uváděné pájky SAC305 a SN100C.
\end{itemize}
\textbf{Složky tavidla:} nosič (přírodní nebo syntetická pryskyřice, případně báze organických kyselin), aktivátory (organické kyseliny, aminy, halogenidové složky), aditiva a rozpouštědla (např. izopropanol, voda či kombinace).
